\selectlanguage{spanish}

\renewcommand{\articuloidiomaxml}{es}
\renewcommand{\articulotitulo}{Las deudas pendientes de la digitalización en América Latina}
\renewcommand{\articuloautores}{Melisa Ferreyra}
\renewcommand{\articulodoi}{10.1080/10749039.2021.1908364}
\renewcommand{\articulotipo}{EDITORIAL}
\renewcommand{\articulorevista}{Actualidad Latinoamericana}
\renewcommand{\articuloissne}{2618-3137}
\renewcommand{\articulovolumen}{01}
\renewcommand{\articulonumero}{01}
\renewcommand{\articulofecha}{ 2026}
\renewcommand{\articulopagina}{1}

\setcounter{page}{1}
\thispagestyle{firstpage}

% --- TÍTULO A ANCHO COMPLETO ---
\noindent\begin{minipage}[t]{\dimexpr\textwidth+\marginparsep+\marginparwidth\relax}%
{\sffamily\bfseries\Large\articulotitulo\par}%
\end{minipage}\par
\vspace{3mm}%
\renewcommand{\articulotitulotrans}{The Unresolved Challenges of Digitalization in Latin America}
\noindent\begin{minipage}[t]{\dimexpr\textwidth+\marginparsep+\marginparwidth\relax}%
{\sffamily\itshape\normalsize\articulotitulotrans\par}%
\end{minipage}\par
\vspace{4mm}%

% --- BLOQUE LATERAL DE METADATOS ---
\begin{textblock*}{68mm}(\gbbloquex,92mm)
{\setlength{\leftskip}{0pt}\setlength{\parindent}{0pt}%
\noindent\begin{minipage}[t]{68mm}
\sffamily\small\setlength{\parskip}{1pt}
{\bfseries\color{azulrevista}Melisa Ferreyra}\par
{\scriptsize\texttt{\href{https://orcid.org/0000-0002-1825-0097}{https://orcid.org/0000-0002-1825-0097}}}\par
Edición, Universidad de Buenos Aires, Buenos Aires, Argentina\par
\vspace{6pt}
{\bfseries\color{azulrevista!70}Derechos de autor\'ia:}\par
\textcopyright\ 2026 Ferreyra, M. Publicado por Actualidad Latinoamericana. Se permite el uso, distribución y reproducción sin restricciones, incluso con fines comerciales, siempre que se cite la obra original. (\href{https://creativecommons.org/licenses/by/4.0/}{https://creativecommons.org/licenses/by/4.0/})\par\vspace{6pt}
{\bfseries\color{azulrevista!70}Publicación:} 2026, vol.~01, n\'um.~01, pp.~1--4\par\vspace{6pt}
{\bfseries\color{azulrevista!70}URL:} {\scriptsize\texttt{\href{https://melcferreyra.github.io/r-01-algo-v01-n01/}{https://melcferreyra.github.io/r-01-algo-v01-n01/}}}\par\vspace{6pt}
{\bfseries\color{azulrevista!70}DOI:} {\scriptsize\texttt{\href{https://doi.org/10.1080/10749039.2021.1908364}{10.1080/10749039.2021.1908364}}}\par\vspace{6pt}
{\bfseries\color{azulrevista!70}Licencia:} CC by 4.0\par
{\scriptsize\texttt{\href{https://creativecommons.org/licenses/by/4.0/}{https://creativecommons.org/licenses/by/4.0/}}}\par
\vspace{6pt}
{\bfseries\color{azulrevista!70}Fechas}\par\vspace{2pt}
Recibido: 05/06/2026\par
Aceptado: 09/06/2026\par
Online: 26/06/2026\par
\vspace{6pt}
\vspace{6pt}
{\bfseries\color{azulrevista!70}Compartir:}\par\vspace{3pt}
{\large
\href{https://bsky.app/intent/compose?text=https%3A%2F%2Fmelcferreyra.github.io%2Fr-01-algo-v01-n01%2F}{\includesvg[height=0.9em]{/home/melisa/.gbpublisher/fonts/bluesky}}~
\href{mailto:?subject=Las%20deudas%20pendientes%20de%20la%20digitalizaci%C3%B3n%20en%20Am%C3%A9rica%20Latina&body=https%3A%2F%2Fmelcferreyra.github.io%2Fr-01-algo-v01-n01%2F}{\includesvg[height=0.9em]{/home/melisa/.gbpublisher/fonts/envelope}}~
\href{https://www.facebook.com/sharer/sharer.php?u=https%3A%2F%2Fmelcferreyra.github.io%2Fr-01-algo-v01-n01%2F}{\includesvg[height=0.9em]{/home/melisa/.gbpublisher/fonts/facebook}}~
\href{https://www.linkedin.com/sharing/share-offsite/?url=https%3A%2F%2Fmelcferreyra.github.io%2Fr-01-algo-v01-n01%2F}{\includesvg[height=0.9em]{/home/melisa/.gbpublisher/fonts/linkedin}}~
\href{https://www.reddit.com/submit?url=https%3A%2F%2Fmelcferreyra.github.io%2Fr-01-algo-v01-n01%2F&title=Las%20deudas%20pendientes%20de%20la%20digitalizaci%C3%B3n%20en%20Am%C3%A9rica%20Latina}{\includesvg[height=0.9em]{/home/melisa/.gbpublisher/fonts/reddit}}~
\href{https://www.researchgate.net/search?q=Las%20deudas%20pendientes%20de%20la%20digitalizaci%C3%B3n%20en%20Am%C3%A9rica%20Latina}{\includesvg[height=0.9em]{/home/melisa/.gbpublisher/fonts/researchgate}}~
\href{https://t.me/share/url?url=https%3A%2F%2Fmelcferreyra.github.io%2Fr-01-algo-v01-n01%2F&text=Las%20deudas%20pendientes%20de%20la%20digitalizaci%C3%B3n%20en%20Am%C3%A9rica%20Latina}{\includesvg[height=0.9em]{/home/melisa/.gbpublisher/fonts/telegram}}~
\href{https://wa.me/?text=https%3A%2F%2Fmelcferreyra.github.io%2Fr-01-algo-v01-n01%2F}{\includesvg[height=0.9em]{/home/melisa/.gbpublisher/fonts/whatsapp}}~
\href{https://twitter.com/intent/tweet?url=https%3A%2F%2Fmelcferreyra.github.io%2Fr-01-algo-v01-n01%2F&text=Las%20deudas%20pendientes%20de%20la%20digitalizaci%C3%B3n%20en%20Am%C3%A9rica%20Latina}{\includesvg[height=0.9em]{/home/melisa/.gbpublisher/fonts/x-twitter}}
}\par
\end{minipage}%
}% FIN GRUPO leftskip
\end{textblock*}


\begin{tcolorbox}[colback=brown!8!white,colframe=brown!8!white,boxrule=0pt,arc=2pt,left=4pt,right=4pt,top=3pt,bottom=3pt,before skip=4pt,after skip=4pt]
\begin{otherlanguage}{spanish}
\sffamily\small \textbf{Resumen:} La pandemia de COVID-19 aceleró la digitalización en América Latina y evidenció las desigualdades existentes en el acceso, uso y apropiación de tecnologías digitales. Este editorial analiza cómo la expansión de la conectividad no logró eliminar brechas vinculadas a ingresos, territorio, competencias digitales y disponibilidad de dispositivos. Se plantea que la transformación digital requiere políticas públicas integrales que articulen infraestructura, formación, capacidades institucionales y contenidos pertinentes, evitando reproducir desigualdades preexistentes. Los artículos reunidos en este número abordan estas tensiones y contribuyen al debate sobre los desafíos pendientes de la inclusión digital en la región.
\par\smallskip\noindent\textbf{Palabras clave:} transformación digital, brecha digital, América Latina, desigualdad, inclusión digital
\end{otherlanguage}
\end{tcolorbox}

\begin{tcolorbox}[colback=brown!8!white,colframe=brown!8!white,boxrule=0pt,arc=2pt,left=4pt,right=4pt,top=3pt,bottom=3pt,before skip=4pt,after skip=4pt]
\begin{otherlanguage}{english}
\sffamily\small \textbf{Abstract:} The COVID-19 pandemic accelerated digitalization in Latin America and exposed existing inequalities in access to, use of, and engagement with digital technologies. This editorial examines how the expansion of connectivity has not eliminated gaps related to income, territory, digital skills, and device availability. It argues that digital transformation requires comprehensive public policies that integrate infrastructure, training, institutional capacities, and relevant content, avoiding the reproduction of pre-existing inequalities. The articles gathered in this issue address these tensions and contribute to the debate on the ongoing challenges of digital inclusion in the region.
\par\smallskip\noindent\textbf{Keywords:} digital transformation, digital divide, Latin America, inequality, digital inclusion
\end{otherlanguage}
\end{tcolorbox}

\bigskip
\noindent\rule{\linewidth}{0.4pt}
\bigskip

La pandemia de COVID-19 funcionó como un acelerador involuntario de procesos que llevaban años en gestación. En cuestión de semanas, sistemas educativos enteros migraron a plataformas virtuales, trámites administrativos que exigían presencialidad se trasladaron a portales web y la telemedicina dejó de ser una promesa para convertirse en una necesidad imperiosa. Sin embargo, esa aceleración no hizo más que volver más visibles las fracturas que ya existían en el tejido social de la región.

América Latina llegó a la pandemia con una infraestructura digital profundamente desigual. Mientras las capitales y las grandes áreas metropolitanas disponían de conexiones de banda ancha razonablemente estables, vastas zonas rurales e incluso periferias urbanas carecían de conectividad básica. La brecha no era solo de acceso físico: se manifestaba también en las competencias digitales de la población, en la disponibilidad de dispositivos adecuados y en la capacidad institucional para sostener servicios en línea de manera continua y segura.\autocite{2566-CEPAL2021}

Tres años después del momento más agudo de la crisis sanitaria, el balance resulta ambivalente. Por un lado, se registró una expansión significativa de la conectividad: varios países de la región superaron tasas de penetración de Internet que parecían lejanas a comienzos de 2020. Por otro, las desigualdades de uso —lo que la literatura especializada denomina “brecha digital de segundo nivel”— no solo persistieron sino que, en algunos casos, se profundizaron. Los sectores de menores ingresos accedieron mayoritariamente a través de teléfonos móviles con planes de datos limitados, lo que restringía severamente sus posibilidades de participar en actividades educativas, laborales o de ejercicio ciudadano que demandan conexiones estables y pantallas de mayor tamaño.\footnote{Según estimaciones de la CEPAL, en 2022 el costo del servicio de
banda ancha para la población del primer quintil de ingresos
representaba en promedio el 14\% de su ingreso mensual en la región, una
proporción que vuelve inviable cualquier pretensión de universalidad
efectiva.}

El problema, conviene insistir, no es exclusivamente tecnológico. La digitalización que no va acompañada de políticas públicas integrales —que atiendan simultáneamente la conectividad, la formación, el diseño institucional y la producción de contenidos pertinentes— corre el riesgo de reproducir y amplificar las desigualdades preexistentes. En una región donde la matriz de la desigualdad opera simultáneamente por ingreso, género, etnia, territorio y edad, suponer que la mera provisión de infraestructura resuelve el acceso es una simplificación que ningún responsable de política pública debería permitirse.

Este número de la revista reúne contribuciones que abordan distintas aristas de esta problemática. Los trabajos que aquí se presentan comparten una premisa: la transformación digital de nuestras sociedades no puede analizarse al margen de las condiciones materiales y simbólicas en las que se despliega. Esperamos que su lectura contribuya a enriquecer un debate que, lejos de cerrarse, se vuelve más urgente con cada innovación tecnológica que promete —una vez más— resolver por sí sola las deudas históricas de la región.

